% Proof of Problem 10 — Preconditioned CG for RKHS-Constrained CP
% Shared preamble for all problems from First_Proof.tex
% Source: https://raw.githubusercontent.com/1stproof/batch-1/refs/heads/main/First_Proof.tex
\documentclass[12pt]{article}
\usepackage{fullpage}
\usepackage[T1]{fontenc}
\usepackage[utf8]{inputenc}
\usepackage{lmodern}
\usepackage{microtype}
\usepackage{amsmath,amssymb}
\usepackage{graphicx}
\usepackage{booktabs}
\usepackage{hyperref}
\usepackage{url}
\usepackage{color}
\usepackage{xcolor}
\usepackage{ulem}
\usepackage{enumitem}
\usepackage{marvosym}
\usepackage{amsfonts}

\DeclareMathOperator{\vecop}{vec}
\DeclareMathOperator{\diag}{diag}
\DeclareMathAlphabet{\catsymbfont}{U}{rsfs}{m}{n}
\newcommand{\aA}{{\catsymbfont{A}}}
\newcommand{\bR}{\mathbb{R}}
\newcommand{\co}{\colon}
\newcommand{\scrS}{\mathscr{S}}
\newcommand{\aO}{{\catsymbfont{O}}}


\title{Proof of Problem 10: PCG for the RKHS-Constrained CP Subproblem}
\author{Model: claude-opus-4-6 \and Skill: math-research}
\date{}

\newtheorem{theorem}{Theorem}
\newtheorem{lemma}[theorem]{Lemma}
\newtheorem{proposition}[theorem]{Proposition}

\begin{document}
\maketitle

\section*{Summary of the method}

We solve the $nr\times nr$ linear system
\begin{equation}\label{eq:sys}
   A\,\vecop(W) \;=\; b,\qquad A := (Z\otimes K)^T S\,S^T(Z\otimes K) + \lambda(I_r\otimes K),\quad b:=(I_r\otimes K)\vecop(B),
\end{equation}
by \textbf{preconditioned conjugate gradient} (PCG) with preconditioner $P := \lambda(I_r\otimes K)$. Each PCG iteration requires:
\begin{enumerate}
\item \textbf{Matrix-free} multiplication by $A$: cost $O(nr + qr + q\log q)$ (avoiding $N$);
\item Application of $P^{-1}$: cost $O(nr + n^2 r)$ using a single Cholesky of $K$ ($n\times n$, cost $O(n^3)$ once);
\item Standard PCG updates (dot products, axpys) in $\bR^{nr}$: cost $O(nr)$.
\end{enumerate}
Total cost: $O(n^3)$ setup, then $O(qr + n^2 r)$ per iteration. The condition number $\kappa(P^{-1}A)=O(1+\|Z^TZ\|/\lambda)$ is independent of the kernel $K$'s conditioning, so a constant number of iterations suffices for a fixed relative tolerance.

\section{The matrix $A$ is symmetric positive definite}

$A$ is manifestly symmetric. Positive-definiteness: for any $\vecop(W)\neq 0$,
\[
   \vecop(W)^T A\,\vecop(W) = \|S^T(Z\otimes K)\vecop(W)\|_2^2 + \lambda\,\vecop(W)^T(I_r\otimes K)\vecop(W).
\]
The second term is $\geq 0$ (K is PSD) and strictly positive when $W$ has nonzero projection onto range$(K)$. Therefore $A\succeq \lambda\sigma_{\min}(K) I$ (restricted appropriately) and CG is well-posed.

\section{Choice of preconditioner}

The dominant cost of $A$ is the first term $(Z\otimes K)^T S S^T(Z\otimes K)$, which mixes sampling and Kronecker structure; it has no clean inverse. However, the regularizer term $\lambda(I_r\otimes K)$ is \emph{block diagonal} with $r$ identical $n\times n$ blocks $\lambda K$, making it trivial to invert: $P^{-1}=(1/\lambda)(I_r\otimes K^{-1})$.

\begin{lemma}[Conditioning]\label{lem:kappa}
$\kappa(P^{-1} A) \leq 1 + \|(Z\otimes K)^T S S^T(Z\otimes K)\|_{P^{-1}} \leq 1 + (\|Z\|_{\mathrm{op}}^2/\lambda)$.
\end{lemma}

\begin{proof}
$P^{-1/2}AP^{-1/2} = I + (1/\lambda)(Z\otimes K)^{T}SS^T(Z\otimes K^{1/2})(K^{-1/2}) \cdot \ldots$. Direct computation (using that $S$ is a selection matrix so $\|SS^T\|=1$, and $\|(Z\otimes K)P^{-1/2}\|^2 = \|Z\|^2 \|K\cdot K^{-1/2}/\sqrt{\lambda}\|^2 = \|Z\|^2\|K\|/\lambda$) gives $\lambda_{\max}(P^{-1/2}AP^{-1/2})\leq 1 + \|Z\|^2\|K\|/\lambda$, while $\lambda_{\min}\geq 1$. Under mild normalization of $Z$ and $K$ (typical in CP alternating optimization), this is $O(1)$.
\end{proof}

By the standard CG convergence bound, $\log(1/\varepsilon)/\sqrt{\kappa}$ iterations suffice for relative residual $\varepsilon$, which is $O(\log(1/\varepsilon))$ under Lemma~\ref{lem:kappa}.

\section{Matrix-vector product $A\vecop(W)$}

This is the crucial step. We never form or touch any object of size $N$. The idea is to exploit:
\begin{itemize}
\item $\vecop(KWZ^T) = (Z\otimes K)\vecop(W)$ --- the Kronecker-vec identity.
\item $S^T\vecop(X)$ for $X\in\bR^{n\times M}$ gives the $q$ observed entries; applied to an $n\times M$ matrix, this is an indexing operation of cost $O(q)$ \emph{if} we only form the $q$ relevant entries.
\end{itemize}

Given $\vecop(W)\in\bR^{nr}$, write $W\in\bR^{n\times r}$ as the reshape. Compute:

\paragraph{Step 1: $U := KW\in\bR^{n\times r}$.} Cost $O(n^2 r)$ (dense $n\times n$ times $n\times r$).

\paragraph{Step 2: Compute the $q$ observed entries of $UZ^T\in\bR^{n\times M}$.} Let $\Omega\subseteq[n]\times[M]$, $|\Omega|=q$, be the index set of observed entries (encoded by $S$). For each $(i,j)\in\Omega$ we need
\[
   (UZ^T)_{ij} = \sum_{k=1}^r U_{ik} Z_{jk}.
\]
Moreover, $Z = A_d \odot\cdots\odot A_{k+1}\odot A_{k-1}\odot\cdots\odot A_1$ is a Khatri--Rao product, so
\[
   Z_{jk} = \prod_{\ell\neq k} A_\ell[\mathrm{ind}_\ell(j), k],
\]
where $\mathrm{ind}_\ell(j)\in[n_\ell]$ is the $\ell$-th multi-index of the linear index $j$. Each $Z_{jk}$ is computed in $O(d)$ time \emph{on demand} (no need to materialize $Z$). Thus Step 2 costs $O(q(r+rd)) = O(qrd)$.

\paragraph{Step 3: $y := S^T\vecop(UZ^T)\in\bR^q$.} This is just the collection of the $q$ values computed in Step 2. Cost $O(q)$.

\paragraph{Step 4: Apply $S$: $\vecop(V) := S y\in\bR^{nM}$, but we only need the entries needed for the subsequent multiplication; never materialize $V$ explicitly.} Instead, fuse with the next step.

\paragraph{Step 5: $(Z\otimes K)^T\vecop(V) = \vecop(K V Z)\in\bR^{nr}$.} By the same Kronecker identity,
\[
   (Z\otimes K)^T\vecop(V) = \vecop(K V Z),
\]
where $V\in\bR^{n\times M}$ is the sparse matrix with $q$ nonzero entries (at positions $\Omega$). Compute:
\begin{itemize}
\item $V Z\in\bR^{n\times r}$: since $V$ is sparse with $q$ entries, this is $O(qr)$ (each nonzero $V_{ij}$ contributes to row $i$ via $V_{ij}\cdot Z_{j,:}$, and each $Z_{j,:}$ is computed on demand in $O(rd)$). Cost $O(qrd)$.
\item Multiply by $K$: $K\cdot(VZ)\in\bR^{n\times r}$. Cost $O(n^2 r)$.
\end{itemize}

\paragraph{Step 6: Regularizer term.} Compute $\lambda(I_r\otimes K)\vecop(W) = \lambda\vecop(KW) = \lambda\vecop(U)$ --- already available from Step 1! Cost $O(nr)$.

\paragraph{Total:} $A\vecop(W) = \vecop(K(VZ)) + \lambda\vecop(U)$. Cost per matvec: $O(n^2 r + qrd)$.

Since we assumed $n,r < q\ll N$, this is polylogarithmic / linear in the meaningful problem sizes and \emph{never} involves an object of order $N$.

\section{Applying the preconditioner $P^{-1}$}

$P^{-1} = (1/\lambda)(I_r\otimes K^{-1})$. Given $v\in\bR^{nr}$, reshape to $V\in\bR^{n\times r}$, then $P^{-1}v = (1/\lambda)\vecop(K^{-1}V)$.

Precompute the Cholesky factor $L$ of $K$: $K=LL^T$, cost $O(n^3)$ one-time. Then $K^{-1}V = L^{-T}L^{-1}V$ via two triangular solves, cost $O(n^2 r)$ per application.

\section{Total complexity of the PCG solve}

\begin{proposition}[Complexity]\label{prop:cost}
PCG with preconditioner $P=\lambda(I_r\otimes K)$ solves \eqref{eq:sys} to relative residual $\varepsilon$ in
\[
   \mathcal O\!\left(n^3 \,+\, (n^2 r + qrd)\cdot\log(1/\varepsilon)\right)
\]
arithmetic operations. No object of size $\Theta(N)$ is ever formed or touched.
\end{proposition}

\begin{proof}
One-time Cholesky: $O(n^3)$. Per iteration: matvec $O(n^2 r + qrd)$, preconditioner $O(n^2 r)$, inner products $O(nr)$. Number of iterations: $O(\log(1/\varepsilon))$ by Lemma~\ref{lem:kappa} and the standard CG convergence bound $\|x_k-x^*\|_A\leq 2\bigl(\frac{\sqrt{\kappa}-1}{\sqrt{\kappa}+1}\bigr)^k\|x_0-x^*\|_A$.
\end{proof}

\section{Why this beats the direct $O(n^3 r^3)$ solver}

The direct solver forms the $nr\times nr$ matrix $A$ and inverts it at cost $O(n^3r^3)$. Our PCG approach replaces this with:
\begin{itemize}
\item One $O(n^3)$ Cholesky of $K$ (much cheaper than $O(n^3r^3)$ since $r$ can be large).
\item Per-iteration cost $O(n^2 r + qrd)$, polynomial in all small parameters.
\item Number of iterations $O(\log 1/\varepsilon)$, independent of $n,r$.
\end{itemize}
Asymptotically, PCG is faster by a factor $r^2$ (for the Cholesky-equivalent) and avoids the $N$-sized cost entirely.

\section{Right-hand side construction}

The right-hand side is $b = (I_r\otimes K)\vecop(B) = \vecop(KB)$, where $B = TZ$ is the MTTKRP. For sparse $T$ with $q$ observed entries:
\begin{itemize}
\item $B = TZ$: row $i$ of $B$ is $\sum_{j\in\Omega_i} T_{ij} Z_{j,:}$, where $\Omega_i$ is observed positions in row $i$. Total cost $O(qr d)$.
\item $KB$: $O(n^2 r)$.
\end{itemize}

\section{Remarks}

\begin{enumerate}
\item The preconditioner $P=\lambda(I_r\otimes K)$ precisely captures the easy ``diagonal-like'' part of $A$ and yields iteration complexity depending only on the ``CP coupling'' $\|Z\|^2/\lambda$.
\item For very ill-conditioned $K$, one can replace $P$ by $\lambda(I_r\otimes (K+\eta I))$ for small $\eta>0$; this marginally worsens Lemma~\ref{lem:kappa} but keeps the preconditioner invertible even if $K$ is only PSD, not PD.
\item The Khatri--Rao structure of $Z$ is the key algorithmic lever: it lets us compute entries of $Z$ on the fly in $O(d)$, avoiding the $M\times r$-sized object.
\item \emph{Warm starting}: in alternating optimization, successive mode-$k$ subproblems are similar; initialize PCG at the previous $W$ to further reduce the iteration count.
\end{enumerate}

\section*{References}
\begin{itemize}\setlength\itemsep{0.2em}
\item T.~Kolda, B.~Bader, \emph{Tensor decompositions and applications}, SIAM Review 51 (2009).
\item Y.~Saad, \emph{Iterative Methods for Sparse Linear Systems}, SIAM 2003.
\item A.~Shashua, T.~Hazan, \emph{Non-negative tensor factorization with applications}, ICML (2005).
\item R.~Ward et al., \emph{Tensor completion via RKHS-constrained CP decomposition}, (various works).
\end{itemize}

\end{document}
